% Options for packages loaded elsewhere
\PassOptionsToPackage{unicode}{hyperref}
\PassOptionsToPackage{hyphens}{url}
\documentclass[
]{article}
\usepackage{xcolor}
\usepackage[margin=1in]{geometry}
\usepackage{amsmath,amssymb}
\setcounter{secnumdepth}{-\maxdimen} % remove section numbering
\usepackage{iftex}
\ifPDFTeX
  \usepackage[T1]{fontenc}
  \usepackage[utf8]{inputenc}
  \usepackage{textcomp} % provide euro and other symbols
\else % if luatex or xetex
  \usepackage{unicode-math} % this also loads fontspec
  \defaultfontfeatures{Scale=MatchLowercase}
  \defaultfontfeatures[\rmfamily]{Ligatures=TeX,Scale=1}
\fi
\usepackage{lmodern}
\ifPDFTeX\else
  % xetex/luatex font selection
\fi
% Use upquote if available, for straight quotes in verbatim environments
\IfFileExists{upquote.sty}{\usepackage{upquote}}{}
\IfFileExists{microtype.sty}{% use microtype if available
  \usepackage[]{microtype}
  \UseMicrotypeSet[protrusion]{basicmath} % disable protrusion for tt fonts
}{}
\makeatletter
\@ifundefined{KOMAClassName}{% if non-KOMA class
  \IfFileExists{parskip.sty}{%
    \usepackage{parskip}
  }{% else
    \setlength{\parindent}{0pt}
    \setlength{\parskip}{6pt plus 2pt minus 1pt}}
}{% if KOMA class
  \KOMAoptions{parskip=half}}
\makeatother
\usepackage{color}
\usepackage{fancyvrb}
\newcommand{\VerbBar}{|}
\newcommand{\VERB}{\Verb[commandchars=\\\{\}]}
\DefineVerbatimEnvironment{Highlighting}{Verbatim}{commandchars=\\\{\}}
% Add ',fontsize=\small' for more characters per line
\usepackage{framed}
\definecolor{shadecolor}{RGB}{248,248,248}
\newenvironment{Shaded}{\begin{snugshade}}{\end{snugshade}}
\newcommand{\AlertTok}[1]{\textcolor[rgb]{0.94,0.16,0.16}{#1}}
\newcommand{\AnnotationTok}[1]{\textcolor[rgb]{0.56,0.35,0.01}{\textbf{\textit{#1}}}}
\newcommand{\AttributeTok}[1]{\textcolor[rgb]{0.13,0.29,0.53}{#1}}
\newcommand{\BaseNTok}[1]{\textcolor[rgb]{0.00,0.00,0.81}{#1}}
\newcommand{\BuiltInTok}[1]{#1}
\newcommand{\CharTok}[1]{\textcolor[rgb]{0.31,0.60,0.02}{#1}}
\newcommand{\CommentTok}[1]{\textcolor[rgb]{0.56,0.35,0.01}{\textit{#1}}}
\newcommand{\CommentVarTok}[1]{\textcolor[rgb]{0.56,0.35,0.01}{\textbf{\textit{#1}}}}
\newcommand{\ConstantTok}[1]{\textcolor[rgb]{0.56,0.35,0.01}{#1}}
\newcommand{\ControlFlowTok}[1]{\textcolor[rgb]{0.13,0.29,0.53}{\textbf{#1}}}
\newcommand{\DataTypeTok}[1]{\textcolor[rgb]{0.13,0.29,0.53}{#1}}
\newcommand{\DecValTok}[1]{\textcolor[rgb]{0.00,0.00,0.81}{#1}}
\newcommand{\DocumentationTok}[1]{\textcolor[rgb]{0.56,0.35,0.01}{\textbf{\textit{#1}}}}
\newcommand{\ErrorTok}[1]{\textcolor[rgb]{0.64,0.00,0.00}{\textbf{#1}}}
\newcommand{\ExtensionTok}[1]{#1}
\newcommand{\FloatTok}[1]{\textcolor[rgb]{0.00,0.00,0.81}{#1}}
\newcommand{\FunctionTok}[1]{\textcolor[rgb]{0.13,0.29,0.53}{\textbf{#1}}}
\newcommand{\ImportTok}[1]{#1}
\newcommand{\InformationTok}[1]{\textcolor[rgb]{0.56,0.35,0.01}{\textbf{\textit{#1}}}}
\newcommand{\KeywordTok}[1]{\textcolor[rgb]{0.13,0.29,0.53}{\textbf{#1}}}
\newcommand{\NormalTok}[1]{#1}
\newcommand{\OperatorTok}[1]{\textcolor[rgb]{0.81,0.36,0.00}{\textbf{#1}}}
\newcommand{\OtherTok}[1]{\textcolor[rgb]{0.56,0.35,0.01}{#1}}
\newcommand{\PreprocessorTok}[1]{\textcolor[rgb]{0.56,0.35,0.01}{\textit{#1}}}
\newcommand{\RegionMarkerTok}[1]{#1}
\newcommand{\SpecialCharTok}[1]{\textcolor[rgb]{0.81,0.36,0.00}{\textbf{#1}}}
\newcommand{\SpecialStringTok}[1]{\textcolor[rgb]{0.31,0.60,0.02}{#1}}
\newcommand{\StringTok}[1]{\textcolor[rgb]{0.31,0.60,0.02}{#1}}
\newcommand{\VariableTok}[1]{\textcolor[rgb]{0.00,0.00,0.00}{#1}}
\newcommand{\VerbatimStringTok}[1]{\textcolor[rgb]{0.31,0.60,0.02}{#1}}
\newcommand{\WarningTok}[1]{\textcolor[rgb]{0.56,0.35,0.01}{\textbf{\textit{#1}}}}
\usepackage{graphicx}
\makeatletter
\newsavebox\pandoc@box
\newcommand*\pandocbounded[1]{% scales image to fit in text height/width
  \sbox\pandoc@box{#1}%
  \Gscale@div\@tempa{\textheight}{\dimexpr\ht\pandoc@box+\dp\pandoc@box\relax}%
  \Gscale@div\@tempb{\linewidth}{\wd\pandoc@box}%
  \ifdim\@tempb\p@<\@tempa\p@\let\@tempa\@tempb\fi% select the smaller of both
  \ifdim\@tempa\p@<\p@\scalebox{\@tempa}{\usebox\pandoc@box}%
  \else\usebox{\pandoc@box}%
  \fi%
}
% Set default figure placement to htbp
\def\fps@figure{htbp}
\makeatother
\setlength{\emergencystretch}{3em} % prevent overfull lines
\providecommand{\tightlist}{%
  \setlength{\itemsep}{0pt}\setlength{\parskip}{0pt}}
\usepackage{bookmark}
\IfFileExists{xurl.sty}{\usepackage{xurl}}{} % add URL line breaks if available
\urlstyle{same}
\hypersetup{
  pdftitle={Data Wrangling},
  pdfauthor={Sinka Khadijah Abubakar},
  hidelinks,
  pdfcreator={LaTeX via pandoc}}

\title{Data Wrangling}
\author{Sinka Khadijah Abubakar}
\date{2026-03-08}

\begin{document}
\maketitle

Question 1. Download two .csv files from Canvas called DiversityData.csv
and Metadata.csv, and read them into R using relative file paths.

\begin{Shaded}
\begin{Highlighting}[]
\CommentTok{\# loading reuired packages}
\FunctionTok{library}\NormalTok{(tidyverse)}
\end{Highlighting}
\end{Shaded}

\begin{verbatim}
## Warning: package 'tidyr' was built under R version 4.5.2
\end{verbatim}

\begin{verbatim}
## Warning: package 'readr' was built under R version 4.5.2
\end{verbatim}

\begin{verbatim}
## Warning: package 'purrr' was built under R version 4.5.2
\end{verbatim}

\begin{verbatim}
## Warning: package 'stringr' was built under R version 4.5.2
\end{verbatim}

\begin{verbatim}
## -- Attaching core tidyverse packages ------------------------ tidyverse 2.0.0 --
## v dplyr     1.1.4     v readr     2.1.6
## v forcats   1.0.0     v stringr   1.6.0
## v ggplot2   4.0.2     v tibble    3.3.0
## v lubridate 1.9.4     v tidyr     1.3.2
## v purrr     1.2.1     
## -- Conflicts ------------------------------------------ tidyverse_conflicts() --
## x dplyr::filter() masks stats::filter()
## x dplyr::lag()    masks stats::lag()
## i Use the conflicted package (<http://conflicted.r-lib.org/>) to force all conflicts to become errors
\end{verbatim}

\begin{Shaded}
\begin{Highlighting}[]
\CommentTok{\# Reading csv files}
\NormalTok{diversity }\OtherTok{\textless{}{-}} \FunctionTok{read.csv}\NormalTok{(}\StringTok{"DiversityData.csv"}\NormalTok{)}
\NormalTok{metadata }\OtherTok{\textless{}{-}} \FunctionTok{read.csv}\NormalTok{(}\StringTok{"Metadata.csv"}\NormalTok{)}
\end{Highlighting}
\end{Shaded}

Question 2. Join the two dataframes together by the common column
`Code'. Name the resulting dataframe alpha.

\begin{Shaded}
\begin{Highlighting}[]
\CommentTok{\# Joining the two dataframes together by the common column ‘Code’}
\NormalTok{alpha }\OtherTok{\textless{}{-}} \FunctionTok{left\_join}\NormalTok{(diversity, metadata, }\AttributeTok{by =} \StringTok{"Code"}\NormalTok{)}
\FunctionTok{head}\NormalTok{(alpha)}
\end{Highlighting}
\end{Shaded}

\begin{verbatim}
##     Code  shannon invsimpson   simpson richness Crop Time_Point Replicate
## 1 S01_13 6.624921   210.7279 0.9952545     3319 Soil          0         1
## 2 S02_16 6.612413   206.8666 0.9951660     3079 Soil          0         2
## 3 S03_19 6.660853   213.0184 0.9953056     3935 Soil          0         3
## 4 S04_22 6.660671   204.6908 0.9951146     3922 Soil          0         4
## 5 S05_25 6.610965   200.2552 0.9950064     3196 Soil          0         5
## 6 S06_28 6.650812   199.3211 0.9949830     3481 Soil          0         6
##   Water_Imbibed
## 1            na
## 2            na
## 3            na
## 4            na
## 5            na
## 6            na
\end{verbatim}

Question 3. Calculate Pielou's evenness index: Pielou's evenness is an
ecological parameter calculated by the Shannon diversity index (column
Shannon) divided by the log of the richness column. a. Using mutate,
create a new column to calculate Pielou's evenness index. b. Name the
resulting dataframe alpha\_even.

\begin{Shaded}
\begin{Highlighting}[]
\CommentTok{\# Creating a new dataframe and calculating Pielou’s evenness index}
\NormalTok{alpha\_even }\OtherTok{\textless{}{-}}\NormalTok{ alpha }\SpecialCharTok{\%\textgreater{}\%}
  \FunctionTok{mutate}\NormalTok{(}\AttributeTok{even =}\NormalTok{ shannon }\SpecialCharTok{/} \FunctionTok{log}\NormalTok{(richness))}
\FunctionTok{head}\NormalTok{(alpha\_even)}
\end{Highlighting}
\end{Shaded}

\begin{verbatim}
##     Code  shannon invsimpson   simpson richness Crop Time_Point Replicate
## 1 S01_13 6.624921   210.7279 0.9952545     3319 Soil          0         1
## 2 S02_16 6.612413   206.8666 0.9951660     3079 Soil          0         2
## 3 S03_19 6.660853   213.0184 0.9953056     3935 Soil          0         3
## 4 S04_22 6.660671   204.6908 0.9951146     3922 Soil          0         4
## 5 S05_25 6.610965   200.2552 0.9950064     3196 Soil          0         5
## 6 S06_28 6.650812   199.3211 0.9949830     3481 Soil          0         6
##   Water_Imbibed      even
## 1            na 0.8171431
## 2            na 0.8232216
## 3            na 0.8046776
## 4            na 0.8049774
## 5            na 0.8192376
## 6            na 0.8155427
\end{verbatim}

Question 4. Using tidyverse language of functions and the pipe, use the
summarise function and tell me the mean and standard error evenness
grouped by crop over time. a. Start with the alpha\_even dataframe b.
Group the data: group the data by Crop and Time\_Point. c.~Summarize the
data: Calculate the mean, count, standard deviation, and standard error
for the even variable within each group. d.~Name the resulting dataframe
alpha\_average

\begin{Shaded}
\begin{Highlighting}[]
\NormalTok{alpha\_average }\OtherTok{\textless{}{-}}\NormalTok{ alpha\_even }\SpecialCharTok{\%\textgreater{}\%}
  \FunctionTok{group\_by}\NormalTok{(Crop, Time\_Point) }\SpecialCharTok{\%\textgreater{}\%}    \CommentTok{\# grouping the data by Crop and Time\_Point}
  \FunctionTok{summarise}\NormalTok{(                        }\CommentTok{\# Summarizing the data}
    \AttributeTok{mean.even =} \FunctionTok{mean}\NormalTok{(even),}
    \AttributeTok{n =} \FunctionTok{n}\NormalTok{(),}
    \AttributeTok{sd =} \FunctionTok{sd}\NormalTok{(even),}
    \AttributeTok{se =}\NormalTok{ sd }\SpecialCharTok{/} \FunctionTok{sqrt}\NormalTok{(n)}
\NormalTok{  )}
\end{Highlighting}
\end{Shaded}

\begin{verbatim}
## `summarise()` has grouped output by 'Crop'. You can override using the
## `.groups` argument.
\end{verbatim}

\begin{Shaded}
\begin{Highlighting}[]
\FunctionTok{head}\NormalTok{(alpha\_average)              }\CommentTok{\# plotting the head of the resulting dataframe}
\end{Highlighting}
\end{Shaded}

\begin{verbatim}
## # A tibble: 6 x 6
## # Groups:   Crop [2]
##   Crop   Time_Point mean.even     n      sd      se
##   <chr>       <int>     <dbl> <int>   <dbl>   <dbl>
## 1 Cotton          0     0.820     6 0.00556 0.00227
## 2 Cotton          6     0.805     6 0.00920 0.00376
## 3 Cotton         12     0.767     6 0.0157  0.00640
## 4 Cotton         18     0.755     5 0.0169  0.00755
## 5 Soil            0     0.814     6 0.00765 0.00312
## 6 Soil            6     0.810     6 0.00587 0.00240
\end{verbatim}

Question 5. Calculate the difference between the soybean column, the
soil column, and the difference between the cotton column and the soil
column a. Start with the alpha\_average dataframe b. Select relevant
columns: select the columns Time\_Point, Crop, and mean.even. c.~Reshape
the data: Use the pivot\_wider function to transform the data from long
to wide format, creating new columns for each Crop with values from
mean.even. d.~Calculate differences: Create new columns named
diff.cotton.even and diff.soybean.even by calculating the difference
between Soil and Cotton, and Soil and Soybean, respectively. e. Name the
resulting dataframe alpha\_average2

\begin{Shaded}
\begin{Highlighting}[]
\NormalTok{alpha\_average2 }\OtherTok{\textless{}{-}}\NormalTok{ alpha\_average }\SpecialCharTok{\%\textgreater{}\%}
  \FunctionTok{select}\NormalTok{(Time\_Point, Crop, mean.even) }\SpecialCharTok{\%\textgreater{}\%}   \CommentTok{\# selecting relevant columns}
  \FunctionTok{pivot\_wider}\NormalTok{(}\AttributeTok{names\_from =}\NormalTok{ Crop, }\AttributeTok{values\_from =}\NormalTok{ mean.even) }\SpecialCharTok{\%\textgreater{}\%}   \CommentTok{\# pivoting wider}
  \FunctionTok{mutate}\NormalTok{(                                   }\CommentTok{\# calculating differences}
    \AttributeTok{diff.cotton.even =}\NormalTok{ Soil }\SpecialCharTok{{-}}\NormalTok{ Cotton,}
    \AttributeTok{diff.soybean.even =}\NormalTok{ Soil }\SpecialCharTok{{-}}\NormalTok{ Soybean}
\NormalTok{  )}
\FunctionTok{head}\NormalTok{(alpha\_average2)                }\CommentTok{\# plotting the head of the resulting dataframe}
\end{Highlighting}
\end{Shaded}

\begin{verbatim}
## # A tibble: 4 x 6
##   Time_Point Cotton  Soil Soybean diff.cotton.even diff.soybean.even
##        <int>  <dbl> <dbl>   <dbl>            <dbl>             <dbl>
## 1          0  0.820 0.814   0.822         -0.00602          -0.00740
## 2          6  0.805 0.810   0.764          0.00507           0.0459 
## 3         12  0.767 0.798   0.687          0.0313            0.112  
## 4         18  0.755 0.800   0.716          0.0449            0.0833
\end{verbatim}

Question6. Connecting it to plots a. Start with the alpha\_average2
dataframe b. Select relevant columns: select the columns Time\_Point,
diff.cotton.even, and diff.soybean.even. c.~Reshape the data: Use the
pivot\_longer function to transform the data from wide to long format,
creating a new column named diff that contains the values from
diff.cotton.even and diff.soybean.even. d.~Create the plot: Use ggplot
and geom\_line() with `Time\_Point' on the x-axis, the column `values'
on the y-axis, and different colors for each `diff' category. The column
named `values' come from the pivot\_longer. The resulting plot should
look like the one to the right.

\begin{Shaded}
\begin{Highlighting}[]
\NormalTok{plot\_data }\OtherTok{\textless{}{-}}\NormalTok{ alpha\_average2 }\SpecialCharTok{\%\textgreater{}\%}
  \FunctionTok{select}\NormalTok{(Time\_Point, diff.cotton.even, diff.soybean.even) }\SpecialCharTok{\%\textgreater{}\%}   \CommentTok{\# selecting columns}
  \FunctionTok{pivot\_longer}\NormalTok{(                                 }\CommentTok{\# pivoting longer}
    \FunctionTok{c}\NormalTok{(diff.cotton.even, diff.soybean.even),     }\CommentTok{\# creating new column diff}
    \AttributeTok{names\_to =} \StringTok{"diff"}
\NormalTok{  )}

\FunctionTok{library}\NormalTok{(ggplot2)                         }\CommentTok{\# plotting}
\FunctionTok{ggplot}\NormalTok{(plot\_data, }\FunctionTok{aes}\NormalTok{(}\AttributeTok{x =}\NormalTok{ Time\_Point, }\AttributeTok{y =}\NormalTok{ value, }\AttributeTok{color =}\NormalTok{ diff)) }\SpecialCharTok{+}
  \FunctionTok{geom\_line}\NormalTok{(}\AttributeTok{size =} \FloatTok{1.2}\NormalTok{) }\SpecialCharTok{+}
  \FunctionTok{xlab}\NormalTok{(}\StringTok{"Time (hrs)"}\NormalTok{) }\SpecialCharTok{+}
  \FunctionTok{ylab}\NormalTok{(}\StringTok{"Difference from Soil in Pielou\textquotesingle{}s Evenness"}\NormalTok{) }\SpecialCharTok{+}
  \FunctionTok{theme\_minimal}\NormalTok{()}
\end{Highlighting}
\end{Shaded}

\begin{verbatim}
## Warning: Using `size` aesthetic for lines was deprecated in ggplot2 3.4.0.
## i Please use `linewidth` instead.
## This warning is displayed once per session.
## Call `lifecycle::last_lifecycle_warnings()` to see where this warning was
## generated.
\end{verbatim}

\pandocbounded{\includegraphics[keepaspectratio]{Data-Wrangling_files/figure-latex/unnamed-chunk-6-1.pdf}}

Question 7. Commit and push a gfm .md file to GitHub inside a directory
called Coding Challenge 5. Provide me a link to your github written as a
clickable link in your .pdf or .docx

\textbf{GitHub Repository}

\href{https://github.com/sza0209/PLAP/blob/main/Assignments/Coding\%20Challenge5/Data-Wrangling.md}{View
my Coding Challenge 5 on GitHub}

\end{document}
